% Options for packages loaded elsewhere
% Options for packages loaded elsewhere
\PassOptionsToPackage{unicode}{hyperref}
\PassOptionsToPackage{hyphens}{url}
\PassOptionsToPackage{dvipsnames,svgnames,x11names}{xcolor}
%
\documentclass[
  letterpaper,
  DIV=11,
  numbers=noendperiod]{scrartcl}
\usepackage{xcolor}
\usepackage{amsmath,amssymb}
\setcounter{secnumdepth}{-\maxdimen} % remove section numbering
\usepackage{iftex}
\ifPDFTeX
  \usepackage[T1]{fontenc}
  \usepackage[utf8]{inputenc}
  \usepackage{textcomp} % provide euro and other symbols
\else % if luatex or xetex
  \usepackage{unicode-math} % this also loads fontspec
  \defaultfontfeatures{Scale=MatchLowercase}
  \defaultfontfeatures[\rmfamily]{Ligatures=TeX,Scale=1}
\fi
\usepackage{lmodern}
\ifPDFTeX\else
  % xetex/luatex font selection
\fi
% Use upquote if available, for straight quotes in verbatim environments
\IfFileExists{upquote.sty}{\usepackage{upquote}}{}
\IfFileExists{microtype.sty}{% use microtype if available
  \usepackage[]{microtype}
  \UseMicrotypeSet[protrusion]{basicmath} % disable protrusion for tt fonts
}{}
\makeatletter
\@ifundefined{KOMAClassName}{% if non-KOMA class
  \IfFileExists{parskip.sty}{%
    \usepackage{parskip}
  }{% else
    \setlength{\parindent}{0pt}
    \setlength{\parskip}{6pt plus 2pt minus 1pt}}
}{% if KOMA class
  \KOMAoptions{parskip=half}}
\makeatother
% Make \paragraph and \subparagraph free-standing
\makeatletter
\ifx\paragraph\undefined\else
  \let\oldparagraph\paragraph
  \renewcommand{\paragraph}{
    \@ifstar
      \xxxParagraphStar
      \xxxParagraphNoStar
  }
  \newcommand{\xxxParagraphStar}[1]{\oldparagraph*{#1}\mbox{}}
  \newcommand{\xxxParagraphNoStar}[1]{\oldparagraph{#1}\mbox{}}
\fi
\ifx\subparagraph\undefined\else
  \let\oldsubparagraph\subparagraph
  \renewcommand{\subparagraph}{
    \@ifstar
      \xxxSubParagraphStar
      \xxxSubParagraphNoStar
  }
  \newcommand{\xxxSubParagraphStar}[1]{\oldsubparagraph*{#1}\mbox{}}
  \newcommand{\xxxSubParagraphNoStar}[1]{\oldsubparagraph{#1}\mbox{}}
\fi
\makeatother

\usepackage{color}
\usepackage{fancyvrb}
\newcommand{\VerbBar}{|}
\newcommand{\VERB}{\Verb[commandchars=\\\{\}]}
\DefineVerbatimEnvironment{Highlighting}{Verbatim}{commandchars=\\\{\}}
% Add ',fontsize=\small' for more characters per line
\usepackage{framed}
\definecolor{shadecolor}{RGB}{241,243,245}
\newenvironment{Shaded}{\begin{snugshade}}{\end{snugshade}}
\newcommand{\AlertTok}[1]{\textcolor[rgb]{0.68,0.00,0.00}{#1}}
\newcommand{\AnnotationTok}[1]{\textcolor[rgb]{0.37,0.37,0.37}{#1}}
\newcommand{\AttributeTok}[1]{\textcolor[rgb]{0.40,0.45,0.13}{#1}}
\newcommand{\BaseNTok}[1]{\textcolor[rgb]{0.68,0.00,0.00}{#1}}
\newcommand{\BuiltInTok}[1]{\textcolor[rgb]{0.00,0.23,0.31}{#1}}
\newcommand{\CharTok}[1]{\textcolor[rgb]{0.13,0.47,0.30}{#1}}
\newcommand{\CommentTok}[1]{\textcolor[rgb]{0.37,0.37,0.37}{#1}}
\newcommand{\CommentVarTok}[1]{\textcolor[rgb]{0.37,0.37,0.37}{\textit{#1}}}
\newcommand{\ConstantTok}[1]{\textcolor[rgb]{0.56,0.35,0.01}{#1}}
\newcommand{\ControlFlowTok}[1]{\textcolor[rgb]{0.00,0.23,0.31}{\textbf{#1}}}
\newcommand{\DataTypeTok}[1]{\textcolor[rgb]{0.68,0.00,0.00}{#1}}
\newcommand{\DecValTok}[1]{\textcolor[rgb]{0.68,0.00,0.00}{#1}}
\newcommand{\DocumentationTok}[1]{\textcolor[rgb]{0.37,0.37,0.37}{\textit{#1}}}
\newcommand{\ErrorTok}[1]{\textcolor[rgb]{0.68,0.00,0.00}{#1}}
\newcommand{\ExtensionTok}[1]{\textcolor[rgb]{0.00,0.23,0.31}{#1}}
\newcommand{\FloatTok}[1]{\textcolor[rgb]{0.68,0.00,0.00}{#1}}
\newcommand{\FunctionTok}[1]{\textcolor[rgb]{0.28,0.35,0.67}{#1}}
\newcommand{\ImportTok}[1]{\textcolor[rgb]{0.00,0.46,0.62}{#1}}
\newcommand{\InformationTok}[1]{\textcolor[rgb]{0.37,0.37,0.37}{#1}}
\newcommand{\KeywordTok}[1]{\textcolor[rgb]{0.00,0.23,0.31}{\textbf{#1}}}
\newcommand{\NormalTok}[1]{\textcolor[rgb]{0.00,0.23,0.31}{#1}}
\newcommand{\OperatorTok}[1]{\textcolor[rgb]{0.37,0.37,0.37}{#1}}
\newcommand{\OtherTok}[1]{\textcolor[rgb]{0.00,0.23,0.31}{#1}}
\newcommand{\PreprocessorTok}[1]{\textcolor[rgb]{0.68,0.00,0.00}{#1}}
\newcommand{\RegionMarkerTok}[1]{\textcolor[rgb]{0.00,0.23,0.31}{#1}}
\newcommand{\SpecialCharTok}[1]{\textcolor[rgb]{0.37,0.37,0.37}{#1}}
\newcommand{\SpecialStringTok}[1]{\textcolor[rgb]{0.13,0.47,0.30}{#1}}
\newcommand{\StringTok}[1]{\textcolor[rgb]{0.13,0.47,0.30}{#1}}
\newcommand{\VariableTok}[1]{\textcolor[rgb]{0.07,0.07,0.07}{#1}}
\newcommand{\VerbatimStringTok}[1]{\textcolor[rgb]{0.13,0.47,0.30}{#1}}
\newcommand{\WarningTok}[1]{\textcolor[rgb]{0.37,0.37,0.37}{\textit{#1}}}

\usepackage{longtable,booktabs,array}
\usepackage{calc} % for calculating minipage widths
% Correct order of tables after \paragraph or \subparagraph
\usepackage{etoolbox}
\makeatletter
\patchcmd\longtable{\par}{\if@noskipsec\mbox{}\fi\par}{}{}
\makeatother
% Allow footnotes in longtable head/foot
\IfFileExists{footnotehyper.sty}{\usepackage{footnotehyper}}{\usepackage{footnote}}
\makesavenoteenv{longtable}
\usepackage{graphicx}
\makeatletter
\newsavebox\pandoc@box
\newcommand*\pandocbounded[1]{% scales image to fit in text height/width
  \sbox\pandoc@box{#1}%
  \Gscale@div\@tempa{\textheight}{\dimexpr\ht\pandoc@box+\dp\pandoc@box\relax}%
  \Gscale@div\@tempb{\linewidth}{\wd\pandoc@box}%
  \ifdim\@tempb\p@<\@tempa\p@\let\@tempa\@tempb\fi% select the smaller of both
  \ifdim\@tempa\p@<\p@\scalebox{\@tempa}{\usebox\pandoc@box}%
  \else\usebox{\pandoc@box}%
  \fi%
}
% Set default figure placement to htbp
\def\fps@figure{htbp}
\makeatother





\setlength{\emergencystretch}{3em} % prevent overfull lines

\providecommand{\tightlist}{%
  \setlength{\itemsep}{0pt}\setlength{\parskip}{0pt}}



 


\usepackage{booktabs}
\usepackage{longtable}
\usepackage{array}
\usepackage{multirow}
\usepackage{wrapfig}
\usepackage{float}
\usepackage{colortbl}
\usepackage{pdflscape}
\usepackage{tabu}
\usepackage{threeparttable}
\usepackage{threeparttablex}
\usepackage[normalem]{ulem}
\usepackage{makecell}
\usepackage{xcolor}
\usepackage{fvextra}
\DefineVerbatimEnvironment{Highlighting}{Verbatim}{breaklines,commandchars=\\\{\}}
\DefineVerbatimEnvironment{OutputCode}{Verbatim}{breaklines,commandchars=\\\{\}}
\KOMAoption{captions}{tableheading}
\makeatletter
\@ifpackageloaded{caption}{}{\usepackage{caption}}
\AtBeginDocument{%
\ifdefined\contentsname
  \renewcommand*\contentsname{Table of contents}
\else
  \newcommand\contentsname{Table of contents}
\fi
\ifdefined\listfigurename
  \renewcommand*\listfigurename{List of Figures}
\else
  \newcommand\listfigurename{List of Figures}
\fi
\ifdefined\listtablename
  \renewcommand*\listtablename{List of Tables}
\else
  \newcommand\listtablename{List of Tables}
\fi
\ifdefined\figurename
  \renewcommand*\figurename{Figure}
\else
  \newcommand\figurename{Figure}
\fi
\ifdefined\tablename
  \renewcommand*\tablename{Table}
\else
  \newcommand\tablename{Table}
\fi
}
\@ifpackageloaded{float}{}{\usepackage{float}}
\floatstyle{ruled}
\@ifundefined{c@chapter}{\newfloat{codelisting}{h}{lop}}{\newfloat{codelisting}{h}{lop}[chapter]}
\floatname{codelisting}{Listing}
\newcommand*\listoflistings{\listof{codelisting}{List of Listings}}
\makeatother
\makeatletter
\makeatother
\makeatletter
\@ifpackageloaded{caption}{}{\usepackage{caption}}
\@ifpackageloaded{subcaption}{}{\usepackage{subcaption}}
\makeatother
\usepackage{bookmark}
\IfFileExists{xurl.sty}{\usepackage{xurl}}{} % add URL line breaks if available
\urlstyle{same}
\hypersetup{
  pdftitle={PS3 Chunk},
  pdfauthor={Abhishu Karki},
  colorlinks=true,
  linkcolor={blue},
  filecolor={Maroon},
  citecolor={Blue},
  urlcolor={Blue},
  pdfcreator={LaTeX via pandoc}}


\title{PS3 Chunk}
\author{Abhishu Karki}
\date{}
\begin{document}
\maketitle


The summary\_table() function is designed to make it easy to create
summary tables for data analysis. It is based on a revised version of
the table‑formatting function I developed when I reviewed my midterm
exam. In many research settings, analysts repeatedly calculate basic
statistics such as sample size, means, and standard deviations for
different groups. This function automates that process by allowing the
user to specify a grouping variable and a numeric variable, and then
produces a clean, presentation‑ready table using tidyverse syntax and
kableExtra.

\subsection{Function}\label{function}

\begin{Shaded}
\begin{Highlighting}[]
\NormalTok{penguins }\OtherTok{\textless{}{-}}\NormalTok{ penguins}

\FunctionTok{library}\NormalTok{(tidyverse)}
\FunctionTok{library}\NormalTok{(knitr)}
\FunctionTok{library}\NormalTok{(kableExtra)}

\NormalTok{summary\_table }\OtherTok{\textless{}{-}} \ControlFlowTok{function}\NormalTok{(data, group\_var, numeric\_var,}
                                  \AttributeTok{caption =} \ConstantTok{NULL}\NormalTok{, }\AttributeTok{digits =} \DecValTok{1}\NormalTok{) \{}

\NormalTok{  summary\_df }\OtherTok{\textless{}{-}}\NormalTok{ data }\SpecialCharTok{|\textgreater{}}
    \FunctionTok{filter}\NormalTok{(}
      \SpecialCharTok{!}\FunctionTok{is.na}\NormalTok{(\{\{ group\_var \}\}),}
      \SpecialCharTok{!}\FunctionTok{is.na}\NormalTok{(\{\{ numeric\_var \}\})}
\NormalTok{    ) }\SpecialCharTok{|\textgreater{}}
    \FunctionTok{group\_by}\NormalTok{(\{\{ group\_var \}\}) }\SpecialCharTok{|\textgreater{}}
    \FunctionTok{summarise}\NormalTok{(}
      \AttributeTok{N =} \FunctionTok{n}\NormalTok{(),}
      \AttributeTok{Mean =} \FunctionTok{round}\NormalTok{(}\FunctionTok{mean}\NormalTok{(\{\{ numeric\_var \}\}), digits),}
      \AttributeTok{SD =} \FunctionTok{round}\NormalTok{(}\FunctionTok{sd}\NormalTok{(\{\{ numeric\_var \}\}), digits),}
      \AttributeTok{.groups =} \StringTok{"drop"}
\NormalTok{    ) }\SpecialCharTok{|\textgreater{}}
    \FunctionTok{rename}\NormalTok{(}\AttributeTok{Group =}\NormalTok{ \{\{ group\_var \}\})}

\NormalTok{  summary\_df }\SpecialCharTok{|\textgreater{}}
    \FunctionTok{kbl}\NormalTok{(}
      \AttributeTok{caption =}\NormalTok{ caption,}
      \AttributeTok{align =} \FunctionTok{c}\NormalTok{(}\StringTok{"l"}\NormalTok{, }\StringTok{"c"}\NormalTok{, }\StringTok{"c"}\NormalTok{, }\StringTok{"c"}\NormalTok{),}
      \AttributeTok{booktabs =} \ConstantTok{TRUE}
\NormalTok{    ) }\SpecialCharTok{|\textgreater{}}
    \FunctionTok{kable\_styling}\NormalTok{(}
      \AttributeTok{position =} \StringTok{"center"}\NormalTok{,}
      \AttributeTok{latex\_options =} \StringTok{"hold\_position"}
\NormalTok{    ) }\SpecialCharTok{|\textgreater{}}
    \FunctionTok{row\_spec}\NormalTok{(}\DecValTok{0}\NormalTok{, }\AttributeTok{bold =} \ConstantTok{TRUE}\NormalTok{, }\AttributeTok{background =} \StringTok{"\#2E86AB"}\NormalTok{, }\AttributeTok{color =} \StringTok{"\#F2F0EF"}\NormalTok{) }\SpecialCharTok{|\textgreater{}}
    \FunctionTok{column\_spec}\NormalTok{(}\DecValTok{1}\NormalTok{, }\AttributeTok{bold =} \ConstantTok{TRUE}\NormalTok{)}
\NormalTok{\}}
\end{Highlighting}
\end{Shaded}

\subsection{demonstration}\label{demonstration}

\begin{Shaded}
\begin{Highlighting}[]

\NormalTok{speciestable }\OtherTok{\textless{}{-}} \FunctionTok{summary\_table}\NormalTok{( penguins, species, body\_mass, }\StringTok{"Mean Body Mass of Penguins by Species (grams)"}\NormalTok{ )}

\NormalTok{gendertable }\OtherTok{\textless{}{-}} \FunctionTok{summary\_table}\NormalTok{( penguins, sex, flipper\_len, }\StringTok{"Mean Flipper Length by Sex (millimeters)"}\NormalTok{ )}

\NormalTok{islandtable }\OtherTok{\textless{}{-}}  \FunctionTok{summary\_table}\NormalTok{( penguins, island , bill\_len, }\StringTok{"Mean Bill Length by Island (millimeters)"}\NormalTok{ )}
\end{Highlighting}
\end{Shaded}

\begin{Shaded}
\begin{Highlighting}[]
\NormalTok{speciestable}
\end{Highlighting}
\end{Shaded}

\begin{table}[!h]
\centering
\caption{\label{tab:unnamed-chunk-2}Mean Body Mass of Penguins by Species (grams)}
\centering
\begin{tabular}[t]{>{}lccc}
\toprule
\cellcolor[HTML]{2E86AB}{\textcolor[HTML]{F2F0EF}{\textbf{Group}}} & \cellcolor[HTML]{2E86AB}{\textcolor[HTML]{F2F0EF}{\textbf{N}}} & \cellcolor[HTML]{2E86AB}{\textcolor[HTML]{F2F0EF}{\textbf{Mean}}} & \cellcolor[HTML]{2E86AB}{\textcolor[HTML]{F2F0EF}{\textbf{SD}}}\\
\midrule
\textbf{Adelie} & 151 & 3700.7 & 458.6\\
\textbf{Chinstrap} & 68 & 3733.1 & 384.3\\
\textbf{Gentoo} & 123 & 5076.0 & 504.1\\
\bottomrule
\end{tabular}
\end{table}

\begin{Shaded}
\begin{Highlighting}[]
\NormalTok{gendertable}
\end{Highlighting}
\end{Shaded}

\begin{table}[!h]
\centering
\caption{\label{tab:unnamed-chunk-2}Mean Flipper Length by Sex (millimeters)}
\centering
\begin{tabular}[t]{>{}lccc}
\toprule
\cellcolor[HTML]{2E86AB}{\textcolor[HTML]{F2F0EF}{\textbf{Group}}} & \cellcolor[HTML]{2E86AB}{\textcolor[HTML]{F2F0EF}{\textbf{N}}} & \cellcolor[HTML]{2E86AB}{\textcolor[HTML]{F2F0EF}{\textbf{Mean}}} & \cellcolor[HTML]{2E86AB}{\textcolor[HTML]{F2F0EF}{\textbf{SD}}}\\
\midrule
\textbf{female} & 165 & 197.4 & 12.5\\
\textbf{male} & 168 & 204.5 & 14.5\\
\bottomrule
\end{tabular}
\end{table}

\begin{Shaded}
\begin{Highlighting}[]
\NormalTok{islandtable}
\end{Highlighting}
\end{Shaded}

\begin{table}[!h]
\centering
\caption{\label{tab:unnamed-chunk-2}Mean Bill Length by Island (millimeters)}
\centering
\begin{tabular}[t]{>{}lccc}
\toprule
\cellcolor[HTML]{2E86AB}{\textcolor[HTML]{F2F0EF}{\textbf{Group}}} & \cellcolor[HTML]{2E86AB}{\textcolor[HTML]{F2F0EF}{\textbf{N}}} & \cellcolor[HTML]{2E86AB}{\textcolor[HTML]{F2F0EF}{\textbf{Mean}}} & \cellcolor[HTML]{2E86AB}{\textcolor[HTML]{F2F0EF}{\textbf{SD}}}\\
\midrule
\textbf{Biscoe} & 167 & 45.3 & 4.8\\
\textbf{Dream} & 124 & 44.2 & 6.0\\
\textbf{Torgersen} & 51 & 39.0 & 3.0\\
\bottomrule
\end{tabular}
\end{table}

\section{Thank You}\label{thank-you}




\end{document}
